% ===================================================================
%  اللوحة نمرة ٩ — الجبل. الغابة.
%
%  Traduction, faite en 2026, en arabe egyptien ecrit, sur l'ido de
%  texto/io. Regles de la colonne : voir 10-lawha-01.tex. Deux
%  scenes, deux numerotations : les cles portent c1 et c2.
%
%  LES PETITES BETES DE LA FORET REDONNENT LA REGLE PHONOLOGIQUE DES
%  TABLEAUX 2, 3 ET 4, ET CETTE FOIS SUR LE MOT LE PLUS COURANT :
%
%    تعبان (7)   le serpent   MSA : ثعبان
%    دبّانة (3)   la mouche    MSA : ذبابة
%    حنش         la couleuvre — le serpent inoffensif, mot propre
%
%  Le ث de ثعبان passe a ت, le ذ de ذبابة a د. Ce sont les memes deux
%  glissements qu'on avait vus sur تلاتة, دقن, دراع, تلج et ديل.
%
%  ET LE CHAMPIGNON (14) S'APPELLE عيش الغراب, « le pain du
%  corbeau ». Le mot عيش est celui du tableau 6 — le pain, qui veut
%  dire la vie — et il sert ici a nommer autre chose : c'est la
%  preuve qu'il est vivant, et pas seulement conserve.
%
%  LA CARRIERE (30) EST المحجر, ET C'EST LE SEUL METIER DU LIVRET QUE
%  L'EGYPTE PRATIQUE DEPUIS CINQ MILLE ANS. Le mot n'a rien
%  d'emprunte et n'a jamais eu besoin d'etre refait.
%
%  LE RESTE DE LA MONTAGNE :
%
%    خُرج (40)    la besace
%    فروة (42)   la pelisse
%    مغارة (56)  la grotte
%    خرايب (14)  les ruines
%    حطّاب (33)   le bucheron
%    فحّام (41)   le charbonnier
%    سايح (46)   le touriste
%    تمّ (45)     le cygne
%    فراولة (11) la fraise
%
%  L'ANNEXION ARABE SUIT L'IDO SANS UNE REPRISE, ET LE TABLEAU 9 LE
%  MONTRE MIEUX QUE LES AUTRES : « منشار الحطّاب », la scie du
%  bucheron, « بوق الصيّادين », le cor des chasseurs — le possede
%  d'abord, le possesseur ensuite, exactement comme l'ido. Les memes
%  quatre blocs avaient coute quatre reprises a la colonne cantonaise.
% ===================================================================

%%K t09-apar-1 apar
\VUsaut{4.0mm}
\VUtitre{15.0pt}{60}{اللوحة نمرة ٩}
\VUsaut{3.0mm}

%%K t09-tit-1 sub
\VUcentre{12.0pt}{0}{\textit{المشهد الأول.}}
\VUsaut{3.0mm}

%%K t09-tit-2 sub
\VUpk{11.4pt}{40}{الجبل. --- بلد المياه.}
\VUsaut{3.0mm}

%%K t09-c1-01-1 p
1. --- بقى لي تمن أيام في براتز. ومش قادر أعبّر عن كل الإعجاب اللي حسّيت بيه
وأنا قدام \VUgras{سلسلة جبال} \textsuperscript{(1)} البرانس الرائعة، اللي
\VUgras{قمّتها} \textsuperscript{(2)} باينة على السما الزرقا زي زخرفة عملاقة.

%%K t09-c1-02-1 p
2. --- وما كنتش متصوّر إن \VUgras{قمم} \textsuperscript{(3)} البرانس بتوصل
\VUgras{ارتفاع} كبير كده، وفضلت مبهور قدام \VUgras{الأنهار الجليدية}
\textsuperscript{(5)} الرائعة و\VUgras{رؤوس الجبال} \textsuperscript{(4)} المغطّاة بالتلج
اللي بتتدحرج عليها \VUgras{انهيارات} مرعبة أوي.

%%K t09-tit-3 sub
\VUsaut{3.0mm}
\VUpk{10.2pt}{40}{منظر الجبل.}
\VUsaut{3.0mm}

%%K t09-c1-03-1 p
3. --- ومن تاني يوم وصولي، رحت لـ\VUgras{البركان} \textsuperscript{(6)} القديم
المطفي اللي جنب براتز. وعلى حواف \VUgras{الفوهة} الجافّة والعارية، لسه باين كويس
أثر مرور \VUgras{الحمم} اللي سالت، من كام قرن، على \VUgras{جنب الجبل}
\textsuperscript{(7)}. وقدرت ألاقي، على واحد من \VUgras{المنحدرات}، كام قطعة من الحمم دي.

%%K t09-c1-04-1 p
4. --- وبراتز واقعة على الحدود، و\VUgras{الحصن} \textsuperscript{(8)} اللي بيحمي
\VUgras{الممر} \textsuperscript{(27)} اللي بيفصل فرنسا عن إسبانيا بيفكّر تمامًا
بالقلاع القديمة على شطّ الراين، اللي شكلها كأنها بتترصّد \VUgras{فريسة} من فوق
صخرتها.

%%K t09-c1-05-1 p
5. --- و\VUgras{نسر} \textsuperscript{(9)} ضخم عشّه مش بعيد عن \VUgras{الطابية}
\textsuperscript{(8)} بيشدّ انتباهي كتير.

%%K t09-c1-05-1 p suite
والطير \VUgras{الجارح} ده، اللي \VUgras{منقاره} قوي، كتير بيجيب لصغاره حمل أو جدي
ماسكه بـ\VUgras{مخالبه}. و\VUgras{جناحيه} كبيرين لدرجة إنه يقدر يحوّم مدة طويلة
وهو شايل حمله التقيل.

%%K t09-tit-4 sub
\VUsaut{3.0mm}
\VUpk{10.2pt}{40}{الرحّالة.}
\VUsaut{3.0mm}

%%K t09-c1-06-1 p
6. --- وأحلى نزهة هنا هي الرحلة لـ\VUgras{شلال} \textsuperscript{(10)} «كاو».
و\VUgras{المية الواقعة} بتنزل بدوّامات وبعدين بتختفي كترعة صغيرة حلوة في
\VUgras{غابة شوح} \textsuperscript{(11)} غرب المدينة. وطلعنا من الغابة دي لحدّ
\VUgras{محطة الأرصاد} \textsuperscript{(12)}، ومن فوق \VUgras{المرصد}، شفنا
\VUgras{منظر} المنطقة كله. و\VUgras{المنظر} رائع، و\VUgras{الطبيعة} كأنها بتغدق
على البلد كل الكنوز اللي عندها.

%%K t09-c1-07-1 p
7. --- وعلشان ننزل، ركبنا \VUgras{الفنيكولير} \textsuperscript{(13)} ووقفنا في محطة
صغيرة جنب \VUgras{خرايب} \textsuperscript{(14)} \VUgras{قلعة إقطاعية} قديمة كان
بيسكنها زمان أمرا براتز.

%%K t09-c1-08-1 p
8. --- يا قلعة مسكينة! ترى بتفكّر في إيه وهي بتبصّ تحت على \VUgras{بلد المية}
\VUgras{الجميلة} \textsuperscript{(15)} اللي بقت دلوقت \VUgras{محطة استشفاء} على الموضة؟

%%K t09-c1-08-2 p
و\VUgras{الجو} مريح أوي، و\VUgras{المية المعدنية} مفيدة أوي، لدرجة إن
\VUgras{العلاج} في \VUgras{المصحّة} \textsuperscript{(17)} بقى متعة حقيقية.

%%K t09-tit-5 sub
\VUsaut{3.0mm}
\VUpk{10.2pt}{40}{بلد مياه مريحة.}
\VUsaut{3.0mm}

%%K t09-c1-09-1 p
9. --- وكمان عملوا كل حاجة علشان الإقامة تبقى مريحة. اتبنى \VUgras{كوبري عالي}
\textsuperscript{(16)} كبير علشان السكة الحديد تعدّي؛ و\VUgras{جنينة}
\textsuperscript{(18)} محطة الاستشفاء، اللي بيتعمل فيها \VUgras{حفلات موسيقى}،
اترتّبت بشكل حلو أوي. واتبنت كام \VUgras{«فيلا»} \textsuperscript{(19)} أنيقة حوالين
\VUgras{الكازينو} \textsuperscript{(21)}؛ و\VUgras{طريق مرصوف} \textsuperscript{(22)}
متعهّد كويس بيسهّل المواصلات أوي.

%%K t09-c1-10-1 p
10. --- و\VUgras{مصطبة} \textsuperscript{(23)} بسيطة بتفصل \VUgras{الطريق}
\textsuperscript{(20)} عن \VUgras{الوادي} \textsuperscript{(25)} نفسه، اللي في قاعه
\VUgras{بحيرة} \textsuperscript{(26)} صغيرة حلوة بتغذّي \VUgras{ترعة}
\textsuperscript{(24)} صافية.

%%K t09-c1-11-1 p
11. --- وعلى الناحية التانية من الوادي، في \VUgras{منجم حديد} \textsuperscript{(28)}
بيتستغلّ. ورحت كذا مرة أشوف \VUgras{عمال المنجم} وهما بينزلوا \VUgras{البير}،
وحتى دخلت \VUgras{الطرقة} اللي بيقضّوا فيها نهارهم، وهما بيطلّعوا \VUgras{الخام}
اللي فيه \VUgras{المعدن}.

%%K t09-c1-12-1 p
12. --- واتبنى \VUgras{مسبك} كبير جنب المنجم علشان يستفيد على طول من الخيرات
اللي بتطلع منه. و\VUgras{الفرن العالي} \textsuperscript{(29)}، اللي بيتسخّن
بـ\VUgras{الفحم} اللي كتير هنا، بيخلّي الحصول على \VUgras{الحديد الزهر} سريع
ورخيص.

%%K t09-c1-13-1 p
13. --- ومش بعيد عن المنجم بيحفروا \VUgras{محجر} \textsuperscript{(30)}.
و\VUgras{الحجّار} العجوز مش بيقطع بـ\VUgras{معوله} لا \VUgras{جرانيت}، ولا
\VUgras{بورفير}، ولا حتى \VUgras{حجر رملي}، لكن \VUgras{حجر مربّع} عادي لبناء
البيوت.

%%K t09-c1-14-1 p
14. --- ومن أحلى زينة البلد دي \VUgras{شجر الشوح} \textsuperscript{(31)}. في كل حتة
في قاع \VUgras{وادي ضيّق} \textsuperscript{(32)}، وعلى شطّ \VUgras{سيل}
\textsuperscript{(33)}، أو \VUgras{هاوية} \textsuperscript{(34)} أو جرف، إبره الرفيّع
بيلوّن الدنيا أخضر.

%%K t09-tit-6 sub
\VUsaut{3.0mm}
\VUpk{10.2pt}{40}{المشي.}
\VUsaut{3.0mm}

%%K t09-c1-15-1 p
15. --- وأنا بستغلّ أجازتي علشان \VUgras{أتمشّى بعيد}. و\VUgras{الطرق}
\textsuperscript{(35)} اللي بتتلوّى على الجبل بتخلّيك تقطع، من غير ما تحسب المسافة،
كذا \VUgras{كيلومتر} وكتير كذا \VUgras{عشرة كيلومترات}.

%%K t09-c1-15-2 p
و\VUgras{الأحجار الحدودية} \textsuperscript{(36)} و\VUgras{أعمدة اللافتات}
\textsuperscript{(37)} بتعلّم الطرق، اللي كتير بتقابل عليها \VUgras{ابن جبل}
\textsuperscript{(38)} شاب و\VUgras{خُرج} \textsuperscript{(40)} على جنبه، شايل
\VUgras{مرموط} \textsuperscript{(39)}؛ أو كام \VUgras{غجري} \textsuperscript{(41)}
\VUgras{فروته} \textsuperscript{(42)} متبايِنة بشكل غريب مع \VUgras{برنسه}
\textsuperscript{(43)} القديم المقطّع. وهو بيسوق بسلسلة \VUgras{دبّ} \textsuperscript{(44)}
مروّض.

%%K t09-tit-7 sub
\VUpk{10.2pt}{40}{تسلّق الجبل.}
\VUsaut{3.0mm}

%%K t09-c1-16-1 p
16. --- وتقريبًا كل يوم بروح مع \VUgras{مرشد} \textsuperscript{(48)} أتمرّن على
\VUgras{التسلّق}، وبفتخر أوي لما، و\VUgras{الشنطة} \textsuperscript{(47)} على ضهري،
و\VUgras{عصاية التسلّق} في إيدي، أهجم زي \VUgras{سايح} \textsuperscript{(46)} حقيقي
على \VUgras{الصخور} \textsuperscript{(45)}.

%%K t09-c1-17-1 p
17. --- ومن فوق الجبل، سكّان السهل ما بيبانوش أكبر من النمل؛
و\VUgras{قطيع خرفان} \textsuperscript{(50)} بيرعى في \VUgras{مرعى} \textsuperscript{(49)}
شكله زي لعبة عيال. و\VUgras{شجرة صنوبر} \textsuperscript{(51)} أو حتى \VUgras{بيت خشب}
\textsuperscript{(52)} ما بيبانوش غير نقطة على الأخضر.

%%K t09-c1-18-1 p
18. --- وفي الرحلات دي، كتير بنقابل على \VUgras{الدرب} \textsuperscript{(53)}
\VUgras{صيّاد الشمواه} \textsuperscript{(54)} شايل على كتفه الحيوان اللي قتله حالًا.

%%K t09-c1-19-1 p
19. --- وحطّيت في مجموعة نباتاتي فرع \VUgras{سرخس} \textsuperscript{(55)} يستاهل
النظر، وفرع \VUgras{خُلّنج} \textsuperscript{(57)} وردي حلو، قطفتهم عند مدخل
\VUgras{مغارة} \textsuperscript{(56)} صغيرة في آخر نزهة لي.

%%K t09-tit-8 sub
\VUsaut{3.0mm}
\VUcentre{12.0pt}{0}{\textit{المشهد التاني.}}
\VUsaut{3.0mm}

%%K t09-tit-9 sub
\VUpk{11.4pt}{40}{الغابة. --- الصيد.}
\VUsaut{3.0mm}

%%K t09-c2-01-1 p
1. --- إمبارح قضّيت يوم لذيذ في الغابة. والنشاط اللي بين الحيوانات
و\VUgras{الحشرات} \textsuperscript{(5)} اللي ساكنة فيها شدّ انتباهي أوي.

%%K t09-c2-01-2 p
\VUgras{العنكبوت} \textsuperscript{(1)} اللي بينسج \VUgras{بيته} \textsuperscript{(2)} بين
فرعين علشان يمسك \VUgras{الدبّانة} \textsuperscript{(3)} اللي بتعدّي؛
و\VUgras{الحلزون} \textsuperscript{(4)} اللي بيشيل قوقعته بتعب؛ والخنفسة اللي بتفتّش
على أكلها؛ والنملة الشغّالة اللي بتلفّ بحماس جنب \VUgras{عش النمل}
\textsuperscript{(6)}؛ الحاجات الصغيّرة دي كلها كانت رايحة جاية وبتشتغل بنشاط مش
عادي.

%%K t09-c2-02-1 p
2. --- و\VUgras{تعبان} \textsuperscript{(7)} فكرته أول ما شفته \VUgras{أفعى}، وطلع
مجرّد \VUgras{حنش} عادي، كان مستخبّي تحت جذع شجرة، ومن وقت للتاني بيطلّع راسه
الرفيّع اللي شكله دايمًا مستعدّ يعضّ. وقدّامي \VUgras{بزّاقة} \textsuperscript{(8)}
كبيرة بتزحف بتقل بعد ما أكلت واحدة من \VUgras{الفراولة} \textsuperscript{(11)} اللذيذة
اللي بتطلع كتير هناك.

%%K t09-c2-03-1 p
3. --- والأرض كانت مفروشة بـ\VUgras{ورد الغابة}: \VUgras{زهرة الربيع}
\textsuperscript{(9)}، و\VUgras{الونكة} \textsuperscript{(10)}، وكمان نباتات كتير ما
كنتش أعرفها، مزهّرة بين \VUgras{خصل الحشيش} \textsuperscript{(12)}.

%%K t09-c2-04-1 p
4. --- وفي المنطقة دي، \VUgras{الكمأة} تقريبًا شايعة زي \VUgras{عيش الغراب}
\textsuperscript{(14)}؛ و\VUgras{خنزير صغير} \textsuperscript{(13)} بتاع واحد غفير كان
بينبش بشراهة يشوف هيلاقي منها ولا لأ.

%%K t09-tit-10 sub
\VUsaut{3.0mm}
\VUpk{10.2pt}{40}{الصيد خلسة.}
\VUsaut{3.0mm}

%%K t09-c2-05-1 p
5. --- وحضرت \VUgras{آخر} مشهد ضحّكني أوي. بمساعدة شمّ \VUgras{كلبه القصير
الرجلين} \textsuperscript{(15)}، \VUgras{الغفير} \textsuperscript{(16)} كان لسه ضابط
\VUgras{صيّاد خلسة} \textsuperscript{(22)}، في اللحظة اللي ده كان بيستعدّ فيها يشيل
\VUgras{الطريدة} \textsuperscript{(17)} اللي قتلها.

%%K t09-c2-05-1 p suite
والصيّاد خلسة، وهو مفضّل إنه يسيب صيده على إنه يتمسك، جري وهرب؛ و\VUgras{الأرنب
البرّي} \textsuperscript{(18)}، و\VUgras{الأيّلة} \textsuperscript{(19)}،
و\VUgras{الوعل} \textsuperscript{(20)}، و\VUgras{الخنزير البرّي} \textsuperscript{(21)}
فضلوا على الحشيش وفرّحوا الغفير.

%%K t09-c2-06-1 p
6. --- وتنوّع الشجر اللي في الغابة بيدهش. \VUgras{الزان} \textsuperscript{(23)}
و\VUgras{البتولا} \textsuperscript{(26)}، و\VUgras{الصنوبر} اللي بيطلّع
\VUgras{الراتينج} \textsuperscript{(27)}، و\VUgras{السنديان} \textsuperscript{(29)}، وكمان
كتير غيرهم، فروعهم متشابكة.

%%K t09-c2-06-2 p
و\VUgras{اللبلاب} \textsuperscript{(24)}، اللي متعلّق في جذوع الشجر العالي، بيلوّن
\VUgras{الغابة} \textsuperscript{(25)} بأخضر لمّاع، اللي بيتّفق بشكل حلو مع الأخضر
الأفتح والأشحب بتاع \VUgras{الطحلب} \textsuperscript{(31)} اللي \VUgras{نقّار الخشب}
\textsuperscript{(30)} بيدوّر فيه على الحشرات.

%%K t09-tit-11 sub
\VUsaut{3.0mm}
\VUpk{10.2pt}{40}{منظر الغابة.}
\VUsaut{3.0mm}

%%K t09-c2-07-1 p
7. --- والسنديان العجوز سحرني. \VUgras{جذوره} \textsuperscript{(32)} التخينة الملتوية،
اللي نُصّها طالع من الأرض، بتديله منظر رائع من القوة والصلابة، وأنا بسأل نفسي
إزاي \VUgras{المالك} \textsuperscript{(35)} يرضى يقطعه ويسلّمه لـ\VUgras{منشار}
\textsuperscript{(34)} \VUgras{الحطّاب} \textsuperscript{(33)}. و\VUgras{الجذوع}
\textsuperscript{(36)} المسكينة، مقشّرة من \VUgras{قشرتها} \textsuperscript{(37)} أو
مشقوقة بـ\VUgras{بلطة} \textsuperscript{(38)} \VUgras{عامل اليومية}
\textsuperscript{(39)}، بتحزّنّي.

%%K t09-c2-08-1 p
8. --- وجنبهم، \VUgras{فحّام} \textsuperscript{(41)} عجوز حضّر بـ\VUgras{جاروفه}
\VUgras{كومة} \textsuperscript{(42)} فحم، وعجوزة شايلة \VUgras{حزمة} \textsuperscript{(40)}
\VUgras{حطب يابس} من فروع الشجر المقطوع.

%%K t09-tit-12 sub
\VUsaut{3.0mm}
\VUpk{10.2pt}{40}{المطاردة.}
\VUsaut{3.0mm}

%%K t09-c2-09-1 p
9. --- وكنت عايز أروح أستريّح شوية في \VUgras{بيت الغفير} \textsuperscript{(46)}، بس
لما وصلت جنب \VUgras{البحيرة الصغيرة} \textsuperscript{(43)}، اللي بيعوم عليها
\VUgras{تمّ} \textsuperscript{(45)} حلو،

%%K t09-c2-09-1 p suite
لاحظت إن \VUgras{فيضان} \textsuperscript{(44)} حصل من قريّب حوّل الطريق لمستنقع
حقيقي. واضطرّيت أتحوّل؛ وأنا بعدّي جنب \VUgras{الأرز} \textsuperscript{(49)}،
و\VUgras{الحور} \textsuperscript{(48)}، و\VUgras{الدردار} \textsuperscript{(47)} اللي على
طرف الغابة، شفت طالع من \VUgras{الغيضة} \textsuperscript{(54)} \VUgras{أيّل}
\textsuperscript{(50)} رائع، وراه \VUgras{كلاب صيد} \textsuperscript{(51)} عنيدة، وكمان
حمّسها \VUgras{بوق} \textsuperscript{(53)} \VUgras{الصيّادين على الخيل}
\textsuperscript{(52)}. والحيوان المسكين كان خلاص مافيش فيه حيل. وجه وقع ميّت على
بعد كام خطوة منّي.

%%K t09-c2-10-1 p
10. --- وابن الغفير، اللي دوشة \VUgras{المطاردة} جذبته، خرج من البيت ورجعنا
احنا الاتنين للغابة. وكان شاف \VUgras{عقعق} \textsuperscript{(56)} على فرع سنديانة
عجوزة. وافتكر إن عشّه على الأغلب مستخبّي في \VUgras{الورق} \textsuperscript{(58)}،
وحبّ ياخد العشّ.

%%K t09-c2-10-2 p
وهو خفيف زي \VUgras{السنجاب} \textsuperscript{(61)}، طلع من \VUgras{فرع}
\textsuperscript{(55)} لفرع، بس ما لقاش حاجة وما عملش غير إنه طيّر \VUgras{غراب}
\textsuperscript{(60)} عجوز كان واقف على \VUgras{غصن} \textsuperscript{(57)}.

\VUsaut{4.0mm}
\VUfilet{22mm}
